\title{Lie algebras and their representations}
\author{\textit{Based on a lecture series by} Ian Grojnowski.}
\date{Michaelmas 2026}
\maketitle

\newcommand{\mat}[1]{\left(\begin{smallmatrix} #1 \end{smallmatrix}\right)}
\newcommand{\Mat}[1]{\textnormal{Mat}_{#1 \times #1}}
\newcommand{\adj}{\textnormal{ad}}
\renewcommand{\gg}{\mathfrak{g}}
\newcommand{\bb}{\mathfrak{b}}
\newcommand{\nn}{\mathfrak{n}}
\renewcommand{\gl}{\mathfrak{gl}}
\renewcommand{\sl}{\mathfrak{sl}}
\newcommand{\so}{\mathfrak{so}}
\renewcommand{\sp}{\mathfrak{sp}}
\renewcommand{\sf}{\mathsf}

This is a course about symmetries controlled by the \emph{linear algebraic groups}
\begin{align*}
    \sf{GL}_n(\bbC) &=\{ A\in \Mat{n}(\bbC) \;|\; \det A \neq 0\}, &
    \sf{O}_n(\bbC) &=\{ A\in \Mat{n}(\bbC) \;|\; A^{\sf{T}} \cdot A = \id \}, \\
    \sf{SL}_n(\bbC) &=\{ A\in \Mat{n}(\bbC) \;|\; \det A = 1\}, &
    \sf{SP}_{n}(\bbC) &=\{ A\in \Mat{n}(\bbC) \;|\; \langle x,y \rangle =
    \langle Ax, Ay\rangle \; \forall \, x,y\in \bbC^{n}\}
\end{align*}
(here the \emph{Symplectic group} $\sf{SP}_n(\bbC)$ is defined
when $n$ is even, and the given definition depends on the choice of a
non-degenerate skew-symmetric bilinear form $\langle-,-\rangle\colon
\bbC^n\times \bbC^n\to \bbC$), and five others --- which happen to be groups
equipped with additional structures of a smooth manifold and an algebraic
variety. These structures mean we can use topological and analytic methods to
study the groups; on the other hand the structures also simplify much of the
study of these groups to simple linear algebra. This reduction is accomplished
by associating the group to its tangent space at identity, which comes naturally
equipped with the structure of a \emph{Lie algebra} that shadows the group
operation. 

This course is about Lie algebras and their representation theory over a field of characteristic zero.

\subsection{About} This course on Lie algebras is a part of Part III of the
Mathematics tripos at the University of Cambridge. The lectures were delivered in-person, and were transcribed live. The notes can be found
online at \url{https://pas201.user.srcf.net/documents/2026-lie-algebras.pdf}. 

Errors and corrections should be communicated to by email to
\url{pas201@srcf.net}.

\section{Getting started}
\marginpar{Lecture 1, 08/10/26} 

The Lie algebra associated to an algebraic group $G$ is the tangent space
$\gg=\sf{T}_\id G$ at the identity element. The group structure on $G$ induces
the structure of a \emph{Lie algebra} on $\gg$. We unpack this with an ad hoc definition first. 

\begin{definition}[Ad hoc] A complex \emph{linear algebraic group} $G$ is a
    subvariety of the affine space $\Mat{n}(\bbC)$ of $n\times n$ matrices,
    defined as the zero locus of some polynomials in matrix coefficients,
    equipped with an associative multiplication map $G\times G\to G$ and an
    inverse map $G\to G$ defined by polynomials in matrix coefficients, such
    that the usual axioms of a group are satisfied.
\end{definition}

This is a terrible definition, because it is not \emph{intrinsic}. It relies on
choosing an embedding of the group into a coordinate system. There \emph{is} an
intrinsic definition, under which a linear algebraic group is simply a
group-object in the category of affine varieties. 

\begin{remark}
    It is not obvious that every such intrinsically defined object can be
    embedded into some affine space $\Mat{n}(\bbC)$, but it follows from noting that
    $G$ would then act on the space of functions on $G$ faithfully, and by
    Hilbert's basis theorem there is a finite dimensional faithful
    sub-representation therein.
\end{remark}

\begin{remark}
    Each linear algebraic group has a maximal compact subgroup (unique up to
    conjugacy), and the linear algebraic group is diffeomorphic to the tangent
    bundle over this compact subgroup. Note that the tangent bundle of any Lie
    group is trivial.

    For instance the maximal compact subgroup of $\sf{SL}_n(\bbC)$ is the special
    unitary group $\sf{SU}_n(\bbC)$, diffeomorphic to an $n$-sphere. Likewise
    the maximal compact subgroup inside $\sf{SO}_n(\bbC)$ is $\sf{SO}_n(\bbR)$.
\end{remark}

\paragraph{What is a tangent space?} Let us begin with an example, by
calculating the tangent space to $\sf{SL}_2(\bbC)$ at $\id=\mat{1&0 \\ 0&1}$. 
An element close to $\id$ can be expressed as \[P=\id + \varepsilon \mat{a & b
\\ c &  d} + \varepsilon^2 ( \cdots ),\]
so for this to lie in $\sf{SL}_2(\bbC)$ we must have $\det(P)=1$, which expands
to
\[1+\varepsilon(a+d)+\varepsilon^2(\cdots )=1.\]
To compute the tangent space we wish for $\varepsilon$ to be ``very small'', so
that $\varepsilon^2\approx 0$ and the tangent space is this first-order neighbourhood
\[\sl_2(\bbC) = \{\mat{a&b \\ c&d} \;|\; a+d =0\}.\] 

This construction can be made rigorous by working over the ring
$\bbE=\bbC[\varepsilon]/(\varepsilon^2)$. Recall that if $G$ is a complex linear
algebraic group and $R$ is a $\bbC$-algebra, we define the \emph{set of
$R$-points of $G$} as 
\[
    G(R)=\{
        P\in \Mat{n}(R) \;|\; 
        P\text{ satisfies the polynomial equations defining } 
        G\subset \Mat{n}(\bbC)
    \}
\]
There is a map $\bbE\to \bbC$ sending $\varepsilon\mapsto 0$, this in turn
induces a map $\pi\colon G(\bbE)\to G(\bbC)$. The tangent space $\gg=\sf{T}_\id G$
is then defined as the pre-image $\pi^{-1}(\id)$.


\begin{exercise}
    Show that $\sf{GL}_n(\bbC)$ is a linear algebraic group (\emph{hint: the condition
    $\det P \neq 0$ can be written as the vanishing of a polynomial by embedding
    $\sf{GL}_n(\bbC)$ into $\Mat{(n+1)}(\bbC)$}). Use this to show that the
    associated Lie algebra is $\frak{gl}_n(\bbC) =\Mat{n}(\bbC)$.
\end{exercise}
 
\begin{exercise}
    For an $n\times n$ matrix $X$, show that $\det(\id + \varepsilon X) = 1 +
    \varepsilon \cdot \text{tr}(X)$. Conclude that the Lie algebra $\sl_n(\bbC)$
    associated to $\sf{SL}_n(\bbC)$ is the space of traceless $n\times n$
    matrices. 
\end{exercise}

\begin{example}
    For the orthogonal group over $\bbC$, the associated Lie algebra is
    the space of antisymmetric matrices. Note that all such matrices are
    traceless over any field of characteristic $\neq 2$, thus $\sf{SO}(n)$ has
    the same Lie algebra. We therefore denote it as $\so_n(\bbC)$.
\end{example}

In the above example we have learnt that different groups can have the same Lie
algebra. This shouldn't be surprising --- for instance, all finite groups are
linear algebraic groups, and their Lie algebra is the zero algebra $\mathbf{0}$.
In a similar fashion the orthogonal group is disconnected, with two connected
components diffeomorphic to each other, one of which is $\sf{SO}_n(\bbC)$.

One thing we will learn in due course is that a Lie group, its maximal compact
subgroup, and its Lie algebra all have coincident finite dimensional
representations.
 

\paragraph{The Lie bracket.} What structure does $\gg=\sf{T}_\id G$ have that is
inherited from the group structure on $G$? It is not multiplication. 

What we do instead is to look at the commutator map $G\times G\to G$ sending
$(p,q)\mapsto p^{-1}q^{-1}pq$. This sends $(\id,\id)\mapsto \id$, thus the
derivative of the commutator gives a linear map 
\[
    \gg\times \gg \simeq T_{(\id,\id)}(G\times
    G)\longrightarrow \sf{T}_\id G \simeq \gg.
\]
This is simply the zero map, and is not quite what we want --- we want a
\emph{bilinear} map $\frak{g}\times \frak{g}\to \frak{g}$, so we differentiate
the first and the second variable separately.

To explain the construction, suppose $G$ is $\sf{GL}_n(\bbC)$ or another matrix
group. Let $P=\id + X\varepsilon$, $Q=\id + Y\delta$ where $\varepsilon, \delta$
are small and satisfy $\varepsilon^2=\delta^2=0$ but
$\varepsilon\delta=\delta\varepsilon\neq 0$ (i.e.\ we work over
the ring $\bbk[\varepsilon, \delta]/(\varepsilon^2,\delta^2)$). We then see that
\[
    P^{-1}Q^{-1}PQ=\id + \varepsilon\delta (XY-YX),
\]
so it makes sense to introduce the map
$[\,,\,]\colon \frak{g}\times
\frak{g}\to \frak{g}$ given by $[X,Y]=XY-YX$. This is the \emph{Lie bracket}.

\begin{exercise}
    The above calculation is an extrinsic proof of the fact that the Lie
    bracket $[\,,\,]\colon\gg\times \gg \to \gg$ is \emph{antisymmetric}. To get an intrinsic proof that doesn't rely
    on any embedding into matrix groups, use the fact that
    $(pqp^{-1}q^{-1})^{-1}=qpq^{-1}p^{-1}$ for any group elements $p,q$.
\end{exercise}

The next property of the Lie bracket is the \emph{Jacobi identity}:
\[[[x,y],z]+[[y,z],x] + [[z,x],y]=0 \qquad \text{for all }x,y,z\in \gg.\] 
We will deduce this over multiple exercises, seeing three reasons why it is
true. Next time we will look at a fourth and a fifth explanation. 

\begin{exercise}
    For $G\subseteq \sf{GL}(n,\bbC)$, show by direct
    computation that the Jacobi identity holds.
\end{exercise}

\begin{exercise}
    Marshall Hall discovered the following identity in the 1930s: writing
    $p^q=q^{-1}pq$ and $(p,q)=p^{-1}q^{-1}pq$, Hall showed that in any
    group $G$ with elements $p,q,r$ we have
    \[((p,q^{-1}),r)^q \cdot ((q,r^{-1}),p)^r \cdot ((r,p^{-1}),q)^p=\id.\] 
    Differentiate this to deduce the Jacobi identity.
\end{exercise}

\begin{exercise}
    Deduce the Jacobi identity from associativity of the group law on $G$.
\end{exercise}

\paragraph{Lie algebras.} The two properties above --- antisymmetry and the
Jacobi identity --- define an abstract Lie algebra. Fix a field $\bbk$ of
characteristic $\neq 2,3$.

\begin{definition}
    Let $\gg$ be a $\bbk$-vector space. A \emph{Lie bracket} on $\gg$ is an
    antisymmetric bilinear form $[\,,\,]\colon \gg\times\gg\to \gg$ that
    satisfies the Jacobi identity. We say $\gg$ equipped with a Lie bracket is a
    \emph{Lie algebra}.
\end{definition}

\begin{exercise}
    \marginpar{Lecture 2, 10/10/26}
    Here is another way to define a Lie bracket on a $\bbk$-vector
    space $\gg$ --- it is given by an element $\mu\in \Exter^2(\gg\dual)\otimes \gg$
    (i.e.\ a $\gg$-valued 2-form on $\gg\dual$) which lies in the kernel of the
    induced exterior derivative $d_\mu \colon \Exter^\bullet(\gg\dual)\otimes \gg \to
    \Exter^{\bullet+1}(\gg\dual)\otimes \gg$. Recall that this derivative is
    given by $d_\mu(1)=0$, $d_\mu\omega = \omega \circ \mu$ for $\omega\in
    \gg\dual\otimes \gg$, extended by the graded Leibniz rule to the rest of
    $\Exter^\bullet(\gg\dual)\otimes \gg$. Check the equivalence of the two
    definitions.
\end{exercise}

\begin{example}
    If $V$ is a finite dimensional $\bbk$-vector space, then $\gl(V) =
    \End_\bbk(V)$ is a Lie algebra, isomorphic to $\gl_n(\bbk)$. We omit the
    field from the notation unless necessary, thus simply writing $\gl_n$. 

    Other examples include the Lie algebras 
    \begin{align*}
        \sl_n &= \{X\in\gl_n\;|\; \text{tr}(X)=0\}, \\ 
        \so_n &= \{X\in \gl_n\;|\; X + X^{\sf{T}} = 0\}, \\ 
        \sp_{n} &= \{ X\in \gl_n\;|\; J\, X^{\sf{T}}\, J^{-1} + X = 0\} \qquad
        (n \text{ even, } J\text{ the standard }n\times n\text{ symplectic
        matrix}), \\ 
        \bb_n &= \{X \in \gl_n\;|\; X \text{ upper triangular}\}, \\ 
        \nn_n &= \{X \in \gl_n\;|\; X \text{ strictly upper triangular}\}.
    \end{align*}
\end{example}

\begin{exercise} 
    Check directly that each of the above, equipped with the commutator bracket,
    is a Lie algebra. Which algebraic groups are they the Lie algebras of?
\end{exercise}

\begin{exercise}
    Show that $\{\mat{\ast & \ast \\ \ast & \cdot}\}$ is not a Lie subalgebra of
    $\gl_2$.
\end{exercise}

\begin{example}
    If $V$ is a $\bbk$-vector space, then the zero form $0\in \Exter^2
    (V\dual)\otimes V$ makes $V$ into a trivial Lie algebra. Such Lie algebras
    are termed \emph{Abelian}. 
\end{example}

\begin{exercise}
    Classify one, two, and three-dimensional Lie algebras over $\bbC$.
\end{exercise}

\section{Representations of Lie algebras}

\begin{definition}
    A \emph{representation} of a Lie algebra $\gg$ on a $\bbk$-vector space $V$ is a
    homomorphism of Lie algebras $\rho \colon \gg\to \gl(V)$. In this case we
    say ``$\gg$ acts on $V$''.
\end{definition}

For example, each Lie algebra $\gg$ acts on itself --- given $x\in \gg$, define
the map $\adj_x\colon \gg \to \gg$ as $y\mapsto [x,y]$.

\begin{lemma}
    The map $\gg \to \gl(\gg)$ given by $x\mapsto \adj_x$ is a representation of
    $\gg$. We call this the \emph{adjoint representation}.
    \begin{proof}
        We must show $\adj_{[x,y]} = \adj_x\circ \adj_y - \adj_y\circ \adj_x$.
        But $\adj_{[x,y]}(z)=[[x,y],z]$, while the right-hand side sends $z$ to
        $[x,[y,z]]-[y,[x,z]]$. By antisymmetry, the required result is
        equivalent to the Jacobi identity.
    \end{proof}
\end{lemma}

\begin{definition}
    The \emph{center} of $\gg$ is the kernel of the adjoint action
    $\ker(\adj\colon \gg\to \gl(\gg))$
\end{definition}

Thus a Lie algebra with trivial center acts faithfully on itself. 

\begin{remark}
In fact \emph{every} finite dimensional Lie algebra has a faithful finite
dimensional representation. This is a non-trivial theorem of Ado, and the
construction of the representation is non-obvious (unlike the
corresponding statement for algebraic groups, where you can simply consider the
space of functions on the group). We will look at the proof if time permits.
\end{remark}

\begin{example}
The Lie algebra $\sl_2=\{\mat{a & b \\ c & -a}\;|\; a,b,c\in \bbk\}$ has a basis given
by the three elements \[e=\mat{0&1\\0&0}, \quad h = \mat{1 & 0 \\ 0 & -1}, \quad f
= \mat{0 & 0 \\ 1 & 0}\]
which satisfy the relations 
\[[h,e]=2e, \quad [h,f]=-2f, \quad [e,f]=h.\]
Hence a representation of $\sl_2$ on $\bbk^n$ is given by a triple $E,F,H\in
\Mat{n}(\bbk)$ such that $EF-FE=H$, $HF-FH=-2F$, and $HE-EH=2E$.
\end{example}

\begin{exercise}
    If $G$ is an algebraic group, an (algebraic) representation of $G$ on a
    $\bbk$-vector space $V$ is a group homomorphism $\rho\colon G\to
    \sf{GL}(V)$ defined by polynomial equations in matrix coefficients. For the
    ring $\bbE=\bbk[\varepsilon]/(\varepsilon^2)$, we thus have a map of the
    groups of $\bbE$-points $G(\bbE)\to \sf{GL}(V)(\bbE)$ satisfying
    $\rho(\id)=\id$. It follows that $\rho(\id + X\varepsilon) = \id  +
    d\rho(X)\cdot \varepsilon$ for some matrix $d\rho(X)$. Check that this
    assignment $d\rho\colon \sf{T}_\id G \to \gl(V)$ is a Lie algebra
    homomorphism.
\end{exercise}

\begin{example}
    Take $G=\sf{SL}_2(\bbk)$, and $L(n)=\Sym^n(\bbk^{\oplus 2})$ (equivalently,
    the degree two homogeneous piece of $\bbk[x,y]$). As $\sf{GL}_2(\bbk)$ acts
    on $L(n)$, so does $\sf{SL}_2(\bbk)$. Call this representation $\rho_n\colon
    \sf{SL}(2,\bbk)\to \sf{GL}(L(n))$. For example, $\rho_0$ is the trivial
    representation, $\rho_1$ is the standard representation, and we have \[\rho_2\mat{a&b \\
    c&d} = \mat{a^2 & ab & b^2 \\ 2ac & ad+bc & 2bd \\ c^2 & cd & d^2}.\]
    Differentiating, we get representations $\sl_n\to \gl(L(n))$ which we
    simply refer to as $L(n)$.
\end{example}

\begin{exercise}
    Writing $(x^iy^j)_{i+j=n}$ for the standard basis of $L(n)$, show that we
    have 
    \begin{gather*}
        e\cdot x^iy^j = jx^{i+1}y^{j-1}, \quad f\cdot x^iy^j = ix^{i-1}y^{j+1},
        \quad h\cdot x^iy^j = (i-j)x^iy^j, \\ 
        \text{or equivalently,}\qquad
        e=x\frac{\partial\phantom{y}}{\partial y}, \quad f = y
        \frac{\partial\phantom{x}}{\partial x}, \quad h =
        x\frac{\partial\phantom{x}}{\partial x} - y
        \frac{\partial\phantom{y}}{\partial y}.
    \end{gather*}
    Assuming $\text{char}(\bbk)=0$, show that each $L(n)$ is an irreducible representation.
\end{exercise}

\begin{remark} 
    How to think about these representations $L(n)$? The space $L(n)$ is
    $\HH^0(\bbP^1, \OO(n))$. This is significant because $\bbP^1$ is a
    homogeneous space for the group $\sf{SL}_2$, and $\sl_2$ is precisely the
    Lie algebra of holomorphic vector fields on $\bbP^1$. A second approach to
    defining Lie algebras and their representations is thus to consider
    homogeneous spaces for other algebraic groups.
\end{remark}

\marginpar{Lecture 3, 13/10/26} 
\begin{example}
    Take the group $G=\bbC^\times$ with its lie algebra $\gg\simeq \bbC$. A
    representation of $\gg$ on a vector space $V$ is an element $A\in \End(V)$,
    and two such elements $A,A'$ define isomorphic representations if and only
    they are conjugate. Thus the isoclasses of $\gg$-representations are in
    bijection with Jordan normal forms of matrices. Likewise a subspace
    $W\subset V$ is a subrepresentation if and only if it is preserved by $A$,
    and since any matrix always has an eigenvector, the irreducible (i.e.\ simple)
    $\gg$-representations are all one-dimensional. In contrast, irreducible
    representations of $G$ are all one-dimensional, but of the form
    $\rho_n\colon z\mapsto z^n$, which have derivative $d\rho_n = n$. Thus out
    of the $\bbC$-many irreducible representations of $\gg$, only $\bbZ$-many
    come from the Lie group $G$.
\end{example}

The situation for $\sl_2$ is drastically different.

\begin{theorem}[(Lie)]
    If $G$ is a simply connected (almost) simple algebraic group over $\bbC$ (or
    a compact Lie group), with Lie algebra $\gg$, then the natural map
    $\textnormal{Rep}(G)\to \textnormal{Rep}(\gg)$ is an equivalence of categories.
\end{theorem}

Here, \emph{almost simple} means that the group has no nontrivial connected
normal subgroup. For example, $\sf{SL}_n(\bbC)$ has centre $\{\pm \id\}$. This
is an important constraint because quotienting by disconnected discrete normal
subgroups does not affect the Lie algebra. The quotient group
$\sf{SL}_n(\bbC)/\{\pm \id\}\simeq \mathsf{PGL}_n(\bbC)$ has the same Lie
algebra as $\sl_2$. This is consistent with the fact that every
$\sf{PGL}_n$-representation is naturally a $\sf{SL}_n$-representation (but this
doesn't account for all $\sf{SL}_n$-representations).

\begin{remark}
    Not all finite dimensional Lie algebras come from Lie groups. In
    characteristic zero, there is a five-dimensional example somewhere in
    Bourbaki. In characteristic $p>0$ it is simpler to see --- noting that the
    $p\times p$ identity matrix is traceless, the quotient $\sl_p/(\bbk\cdot
    \text{id})$ can be shown to not arise from a Lie group.
\end{remark}

\subsection{Classifying representations of $\sl_2(\bbC)$}
 
Let $e,f,h$ be the standard basis of $\sl_2(\bbC)$, with commutation relations
$[e,f]=h$, $[h,e]=2e$, $[h,f]=-2f$.

\begin{theorem}
    Every finite dimensional representation of $\sl_2(\bbC)$ is a direct sum of
    irreducible representations. For all $n\geq 0$, there exists a unique
    irreducible representation of $\sl_2$ that is $(n+1)$-dimensional.
\end{theorem}

We will build up to the proof, by first making some definitions. Thus let $V$ be a
finite dimensional representation of $\sl_2$.

\begin{definition}
    For any complex number $\lambda$, the \emph{$\lambda$-weight space} for $V$
    (denoted $V_\lambda$) is defined as the $\lambda$-eigenspace for the
    operator $h\colon V\to V$.
\end{definition}

Suppose $v\in V_\lambda$. Then the commutation relations imply $ev\in
V_{\lambda+2}$ and $fv\in V_{\lambda-2}$. Thus the actions of $e$ and $f$ move
the vectors `up and down' the $h$-eigenspaces.

\begin{definition}
    If a non-zero vector $v$ lies in $V_\lambda\cap \ker(e)$, i.e.\ satisfies
    $hv=\lambda v$ and $ev = 0$, then we say that $v$ is a \emph{highest weight
    vector of weight $\lambda$}.
\end{definition}

\begin{lemma}
    Let $v\in V_\lambda$ be a highest weight vector of weight $\lambda$. Let
    $W=\langle v, fv, f^2v, \ldots \rangle\subseteq V$. Then $W$ is a
    subrepresentation of $V$.
    \begin{proof}
        We need to show $fW\subseteq W$, $eW\subseteq W$, and $hW\subseteq W$.
        The first and third statements are immediate. To show $eW\subseteq W$,
        show by induction that we have $ef^n v = n(\lambda-n +1)f^{n-1}v$, from
        which the conclusion will follow.
    \end{proof}
\end{lemma}

\begin{lemma}
    Let $v\in V_\lambda$ be a highest weight vector
    of weight $\lambda$. Then $\lambda \in \bbZ_{\geq 0}$ is a non-negative integer.
    \begin{proof}
        All $f^k v$ ($k\geq 0$) lie in distinct eigenspaces of $h$, so that they
        are linearly independent. Since $V$ is finite dimensional, there is a
        $k$ with $f^kv \neq 0$ but $f^{k+1}v= 0$. But then we have \[0=ef^{k+1}v
        = (k+1)(\lambda-k)f^k v,\]
        and since $(k+1)(f^k v)\neq 0$, we must have $\lambda=k$.
    \end{proof}
\end{lemma}

\begin{lemma}
    Every finite dimensional representation of $\sl_2(\bbC)$ has a highest
    weight vector.
    \begin{proof}
        Since $V$ is a $\bbC$-vector space, the operator $h$ has an eigenvector
        $v$. Then $v, ev, e^2v, \ldots$ are all vectors which, if non-zero, are
        linearly independent since they lie in distinct eigenspaces of $h$. But
        $V$ is finite dimensional, so there is a $k$ with $e^k v \neq 0$,
        $e^{k+1}v=0$. Then $e^kv$ is a highest weight vector.
    \end{proof}
\end{lemma}

\begin{corollary}
    If $V$ is an irreducible representation of $\sl_2(\bbC)$ of dimension $n+1$
    ($n\geq 0$), then $V$ has a basis $v_0,v_1,\ldots, v_n$ with $hv_i =
    (n-2i)v_i$, $fv_i=v_{i+1}$, and $ev_i=i(n-i+1)v_{i-1}$. In particular there
    is a unique $(n+1)$-dimensional representation of $\sl_2(\bbC)$.
\end{corollary}

\begin{exercise}
    Write an explicit isomorphism between $L(n)$ and the irreducible
    representation of $\sl_2(\bbC)$ constructed above.
\end{exercise}

We move on to proving \emph{complete reducibility} in $\text{Rep}(\sl_2\bbC)$,
i.e.\ showing that every indecomposable representation is irreducible. When
proving analogous results for finite groups or compact Lie groups, the trick is
to take some average over all group elements. We don't have such tools here.

\begin{definition}
    Let $V$ be a representation of $\sl_2(\bbC)$. The associated \emph{Casimir
    operator}, denoted $\Omega\colon V\to V$, is given by
    $ef+fe+\frac{1}{2}h^2$.
\end{definition}

\begin{lemma}
    The operator $\Omega$ is central, i.e.\ it commutes with $e$, $f$ and $h$.
\end{lemma}
